\documentclass[11pt]{article}
\usepackage{amsmath, amssymb}
\usepackage[margin=2.5cm]{geometry}
\usepackage{xcolor}
\title{Resep Kanvas \\ \large Pusaran Lautan Magma \& Air}
\author{math-aug}
\date{}

\begin{document}
\maketitle

\section{Domain Kanvas}

Kanvas $\mathcal{C}$ adalah ruang piksel persegi panjang pada bidang
$\mathbb{R}^2$:

\begin{equation}
\mathcal{C} = \bigl\{ (x, y) \in \mathbb{R}^2 \;\big|\; 0 \le x < W,\; 0 \le y < H \bigr\}
\end{equation}

dengan
\begin{align}
W &= 7680 \quad \text{(lebar, piksel)} \\
H &= 4320 \quad \text{(tinggi, piksel)}
\end{align}

Origin $(0,0)$ di kiri-atas. Sumbu $x$ ke kanan, sumbu $y$ ke bawah.

\section{Pembagian Tiga Lapisan}

Kanvas terbagi menjadi tiga pita horizontal oleh dua garis horizon
$y_1$ dan $y_2$:

\begin{equation}
y_1 = 1000, \qquad y_2 = 3320
\end{equation}

\begin{align}
\text{Laut Air}   &:\quad \mathcal{L}_{\text{air}}   = \bigl\{ (x,y) \in \mathcal{C} \;\big|\; 0 \le y < y_1 \bigr\} \\
\text{Langit}     &:\quad \mathcal{L}_{\text{langit}}= \bigl\{ (x,y) \in \mathcal{C} \;\big|\; y_1 \le y < y_2 \bigr\} \\
\text{Laut Magma} &:\quad \mathcal{L}_{\text{magma}} = \bigl\{ (x,y) \in \mathcal{C} \;\big|\; y_2 \le y < H \bigr\}
\end{align}

Tinggi tiap lapisan:
\begin{align}
h_{\text{air}}   &= y_1 - 0     = 1000 \\
h_{\text{langit}}&= y_2 - y_1   = 2320 \\
h_{\text{magma}} &= H   - y_2   = 1000
\end{align}

\section{Titik Pusat Pusaran}

Pusat pusaran air $\mathbf{c}_{\text{air}}$ dan pusaran magma
$\mathbf{c}_{\text{magma}}$ masing-masing berada di titik tengah
lapisannya terhadap sumbu $x$:

\begin{align}
\mathbf{c}_{\text{air}}   &= \bigl( \tfrac{W}{2},\; \tfrac{y_1}{2} \bigr)
                            = (3840,\; 500) \\
\mathbf{c}_{\text{magma}} &= \bigl( \tfrac{W}{2},\; \tfrac{y_2 + H}{2} \bigr)
                            = (3840,\; 3820)
\end{align}

Radius maksimum pusaran:
\begin{equation}
r_{\max} = 450
\end{equation}

\section{Simetri Refleksi}

Kedua pusaran simetris terhadap titik tengah kanvas
$\mathbf{m} = \bigl( \tfrac{W}{2}, \tfrac{H}{2} \bigr) = (3840, 2160)$:

\begin{equation}
\mathbf{c}_{\text{magma}} = 2\mathbf{m} - \mathbf{c}_{\text{air}}
\end{equation}

Verifikasi:
\begin{equation}
2(3840, 2160) - (3840, 500) = (3840,\; 3820) = \mathbf{c}_{\text{magma}} \quad \checkmark
\end{equation}

\end{document}
